\section{Proof of the diagonal bound}
\label{sec:concrete-bound}

The local certificate uses the exact parameters
\[
 \begin{aligned}
 p&=\frac{499999999}{10^9},&
 \mu&=\frac{405108059767}{10^{12}},&
 w&=\frac{104160447497}{250000000000},\\
 \beta&=\frac{9999999}{10^7},&
 x&=\frac{337503}{10^6},&
 \xi&=x-\frac4{10^{11}}.
 \end{aligned}
\]
The terminal source parameters, envelope heights, and route data are
recorded in the certificates linked in Section~\ref{sec:formalization}.

\subsection{Certifying the local envelope}

Set $a=2681/10000$. Rational interval bounds give
\[
 \mathcal F(a)>0>\mathcal F(\mu).
\]
Thus $\mu\in F_p\subset(a,\mu]$ by Lemma~\ref{lem:feasibility-cutoff}.

Extend $M_p$ by \eqref{eq:optimizer} to
$D=\{q\in(0,1):p e^{\beta A_p(q)}<U(q)\}$, where the logarithm has
positive argument. Put $H(q)=qM_p(q)+wq$. To prove concavity, set
$G_0(q)=q\log p+\beta q(A_p(q)-B_q)$,
$f(z)=pz+(1-p)\log(1-e^z)$, and
$\bar f(z)=f(\max\{z,\log p\})$. Then
\[
 H(q)=qB_q+wq+
 \frac q\beta\left[\bar f\left(\frac{G_0(q)}q\right)-f(\log p)\right].
\]
The function $G_0$ is convex, since it is
$\beta w[q\log q+(1-q)\log(1-q)]$ plus an affine function;
hence $D$ is an interval. Since $f'(z)=(p-e^z)/(1-e^z)$ and
$f''(z)<0$, the function $\bar f$ is concave and nonincreasing.
The perspective $(q,z)\mapsto q\bar f(z/q)$ is concave for
$q>0$, $z<0$, and nonincreasing in $z$. Since $G_0$ is convex,
$q\bar f(G_0(q)/q)$ is concave on $D$.
Together with concavity of $qB_q=wq\log(\mu/q)$, this proves concavity of $H$.

The certified endpoint inequalities $G_0(a),G_0(\mu)<0$ give
$[a,\mu]\subset D$. Divide this interval into $192$ equal cells
with endpoints $q_i$. At each midpoint take an affine upper support
$T_i$ of $H$. Rational heights $h_0,\ldots,h_{192}$ certify
\[
 H(q)\le T_i(q)\quad(q\in D),\qquad
 T_i(q_i)\le h_i,\qquad T_i(q_{i+1})\le h_{i+1}.
\]
Affine interpolation then bounds $H(q)+wqv$ throughout each cell by
the maximum of its two certified endpoint values, for every real $v$.

Put $a_i=wq_i$, $\Delta_i=a_{i+1}-a_i>0$, and
$r_i=(h_i-h_{i+1})/\Delta_i$. The logarithmic hinge
\[
 E(u)=a_0\log u+h_0+
 \sum_{i=0}^{191}\Delta_i(\log u-r_i)_+
 \qquad(0<u\le1)
\]
dominates every line $a_i\log u+h_i$: telescope the successive line
differences and replace each difference by its positive part.
It is continuous, nondecreasing, and integrable. Integrating a single
hinge gives $\int_0^1(\log u-r)_+\,du\le e^r-r-1$, with equality
for $r\le0$. Telescoping again, the checked exponential sum yields
\begin{equation}\label{eq:certified-integral}
 \int_0^1E(u)\,du
 \le h_{192}-a_{192}+\sum_{i=0}^{191}\Delta_i e^{r_i}
 <-\frac{15}{10^8}.
\end{equation}
Thus $E$ satisfies the hypotheses of Proposition~\ref{prop:local-alternative}.
The exact interval calculations also verify the real-function identities
and the domains of logarithms and inverses.

\subsection{A clique bound at density \texorpdfstring{$1/2$}{1/2}}

Put $z=625145177/500000000$.
\begin{proposition}\label{prop:dense-clique-bound}
There is a constant $C\ge1$ such that, for all sufficiently large $k$,
\[
 d_R(G)\ge\frac12,\quad |V(G)|\ge Ce^{zk}
 \quad\Longrightarrow\quad
 G\text{ contains a red or blue }K_k.
\]
\end{proposition}
\begin{proof}

We reassemble the profile data of~\cite{ZBS}, using Lu--Wang's uniform
baseline~\cite[Lemma~B.1]{LW}. The certificates check the profile
inequalities throughout their domains and the source comparisons for
every target ratio $\rho>0$, with the required strict margins.
The certified family includes four new child profiles and a root,
all satisfying the route conditions.
The root has density $1/2$, a ratio interval containing $1$, and constant
height $z$. Proposition~\ref{prop:joint-construction}, applied to the
enlarged family, establishes its assertion at ratio $1$.
\end{proof}

\begin{proof}[Proof of Theorem~\ref{thm:main}]
One colour of a complete graph has density at least $1/2$.
Proposition~\ref{prop:dense-clique-bound} and $\gamma=z+10^{-9}$ give,
for all sufficiently large $k$,
\[
 R(k,k)\le\left\lceil(C+2)e^{zk}\right\rceil\le e^{\gamma k}.
\]
The exact exponential comparison is
\[
 e^\gamma<3.49135655.
\]
Hence $R(k,k)\le3.49135655^k$ for all sufficiently large $k$.
\end{proof}
